\subsubsection{Normal rings}
\label{editorial-source-section-normal-rings}

\noindent
We first introduce the notion of a normal domain, and then we
introduce the (very general) notion of a normal ring.

\begin{definition}
\label{editorial-source-definition-domain-normal}
A domain $R$ is called {\it normal} if it is integrally
closed in its field of fractions.
\end{definition}

\begin{lemma}
\label{editorial-source-lemma-integral-closure-in-normal}
Let $R \to S$ be a ring map.
If $S$ is a normal domain, then the integral closure of $R$
in $S$ is a normal domain. It also equals the integral closure
of $R$ in the fraction field of the original $S$.
\end{lemma}

\begin{proof}
Write $B\subset S$ for that integral closure, retaining the
original map $R\to S$. The ring $B$ is a domain, since it is
a subring of the domain $S$. The inclusion induces the injective
map $\operatorname{Frac}(B)\to\operatorname{Frac}(S)$,
$a/b\mapsto a/b$; denominators remain nonzero and equality of
fractions is the same equality $ab'=a'b$ in either domain.
If $z\in\operatorname{Frac}(B)$ is integral over $B$, its monic
equation is also a monic equation over $S$. Normality of $S$
therefore places the image of $z$ in $S$. The subalgebra $B[z]$
is integral over $B$, and $R\to B$ is integral by the definition
of $B$. Transitivity (Lemma \ref{algebra-lemma-integral-transitive}) makes
$z$ integral over $R$, so $z\in B$. This proves normality of $B$.
If instead $z\in\operatorname{Frac}(S)$ is integral over $R$,
its same equation, with coefficients mapped into $S$, makes it
integral over $S$. Thus $z\in S$ and then $z\in B$.
The reverse inclusion follows from the original inclusion
$B\subset S\subset\operatorname{Frac}(S)$ and its monic equations,
proving the asserted equality of integral closures.
\end{proof}

\noindent
The following notion is occasionally useful when
studying normality.

\begin{definition}
\label{editorial-source-definition-almost-integral}
Let $R$ be a domain.
\begin{enumerate}
\item An element $g$ of the fraction
field of $R$ is called {\it almost integral over $R$}
if there exists an element $r \in R$, $r\not = 0$
such that $rg^n \in R$ for all $n \geq 0$.
\item The domain $R$ is called {\it completely normal} if every
almost integral element of the fraction field of $R$ is
contained in $R$.
\end{enumerate}
\end{definition}

\noindent
The following lemma shows that a Noetherian domain is normal
if and only if it is completely normal.

\begin{lemma}
\label{editorial-source-lemma-almost-integral}
Let $R$ be a domain with fraction field $K$.
If $u, v \in K$ are almost integral over $R$, then so are
$u + v$ and $uv$. Any element $g \in K$ which is integral over $R$
is almost integral over $R$. If $R$ is Noetherian
then the converse holds as well. More precisely, for an almost
integral $g$ and any original witness $0\ne r\in R$, integrality
of $g$ is equivalent to finite generation of the specified ideal
$I_g=\sum_{n\geq0}R\,rg^n\subset R$.
\end{lemma}

\begin{proof}
Keep witnesses $0\ne r,r'\in R$ with $ru^n\in R$ and
$r'v^n\in R$ for all $n\geq0$. Their product is nonzero, and
the full identities
$$
rr'(uv)^n=(ru^n)(r'v^n),\qquad
rr'(u+v)^n=\sum_{j=0}^n\binom nj(ru^j)(r'v^{n-j})
$$
place both expressions in $R$ for every $n\geq0$. The binomial
coefficients are integers acting through $\mathbf Z\to R$;
no division in $R$ is used. Also $r(-u)^n=(-1)^nru^n\in R$,
and every element of $R$ has witness $1$. Consequently the
almost integral elements form an $R$-subalgebra of the original $K$.
Suppose $g \in K$ is integral over $R$.
In this case there exists a $d > 0$ such that
the ring $R[g]$ is generated by $1, g, \ldots, g^d$ as an $R$-module.
Write each original $g^i=a_i/b_i$ for $0\leq i\leq d$,
with $a_i,b_i\in R$ and $b_i\ne0$, and choose the actual common
denominator $r=\prod_{i=0}^d b_i\ne0$. Then
$rg^i=a_i\prod_{j\ne i}b_j\in R$ for each listed power.
Every element of $R[g]$ is an $R$-linear combination of these
powers, so multiplying that combination by $r$ gives an element
of $R$. In particular $rg^n\in R$ for all $n\geq0$,
which proves that the original $g$ is almost integral.

\medskip\noindent
Let $g\in K$ be almost integral with its original witness
$0\ne r\in R$. The set
$I_g=\sum_{n\geq0}R\,rg^n=rR[g]\subset R$ is an ideal.
Multiplication by $r$ gives the exact $R$-module isomorphism
$$
R[g]\longrightarrow I_g,\quad z\longmapsto rz,
\qquad I_g\longrightarrow R[g],\quad w\longmapsto w/r.
$$
The second map lands in $R[g]$ by the displayed description of
$I_g$, and both composites are identities in the original field $K$.
If $I_g$ is finitely generated, so is $R[g]$, and Lemma
\ref{algebra-lemma-finite-is-integral} proves $g$ integral. Conversely,
if $g$ is integral, Lemma \ref{editorial-source-lemma-characterize-integral}
makes $R[g]$ finite, so the same isomorphism makes $I_g$ finite
for every choice of witness. This proves the exact criterion.
When $R$ is Noetherian the ideal $I_g$ is finitely generated,
giving the claimed converse. Equivalently, the original inclusion
$R[g]\subset(1/r)R$ identifies it with a submodule of the finite
$R$-module $(1/r)R$; multiplication by $r$ identifies this
submodule with the very same ideal $I_g$.
\end{proof}

\begin{lemma}
\label{editorial-source-lemma-localize-normal-domain}
If $R$ is a normal domain and $S\subset R$ is multiplicative
with $0\notin S$, then $S^{-1}R$ is a normal domain.
If $0\in S$, the localization is the zero ring, which is a normal
ring in the sense of Definition \ref{editorial-source-definition-ring-normal}.
\end{lemma}

\begin{proof}
If $0\in S$, every fraction is zero and the spectrum is empty,
so the later definition of a normal ring applies vacuously.
Suppose $0\notin S$ and keep $K=\operatorname{Frac}(R)$.
The original localization injects into $K$ by $a/s\mapsto a/s$.
Its fraction field is identified with $K$ by the inverse maps
$$
\frac{a/s}{b/t}\longmapsto\frac{at}{sb},\qquad
\frac ab\longmapsto\frac{a/1}{b/1}.
$$
All denominators are nonzero; fraction equality in a domain
verifies well-definedness, and multiplication of the original
numerators and denominators verifies both composites.
Let $g\in K$ be integral over $S^{-1}R$, with its monic equation
$P(g)=0$ for $P(x)=x^d+\sum_{0\leq j<d}a_jx^j$, $d\geq1$.
Write $a_j=c_j/s_j$ with $c_j\in R$, $s_j\in S$, and keep
$s=\prod_{j=0}^{d-1}s_j\in S$.
Then $sa_j=c_j\prod_{h\ne j}s_h\in R$, and
$s^{d-j}a_j=s^{d-j-1}(sa_j)\in R$ for every $0\leq j<d$.
The full scaled equation is
$$
(sg)^d+\sum_{j=0}^{d-1}s^{d-j}a_j(sg)^j
=s^d\left(g^d+\sum_{j=0}^{d-1}a_jg^j\right)=0.
$$
It is monic with coefficients in $R$, so normality gives
$sg\in R$. Thus the original $g$ is the image of $(sg)/s$
in $S^{-1}R$, proving the assertion.
\end{proof}

\begin{lemma}
\label{editorial-source-lemma-PID-normal}
A principal ideal domain is normal.
\end{lemma}

\begin{proof}
Let $R$ be a principal ideal domain and retain the original
$g=a/b\in\operatorname{Frac}(R)$, where $b\ne0$, and the full
monic equation
$$
(a/b)^n+\sum_{i=0}^{n-1}r_i(a/b)^i=0,\qquad r_i\in R.
$$
Multiplication by the original $b^n$ gives
$a^n+\sum_{i=0}^{n-1}r_i a^i b^{n-i}=0$ in $R$.
Choose $c$ generating the original ideal $(a,b)$, and write
$a=ca'$, $b=cb'$, $c=ua+vb$ for $a',b',u,v\in R$.
Since $b\ne0$, both $c$ and $b'$ are nonzero. Cancellation in
$c=c(ua'+vb')$ gives $ua'+vb'=1$.
The full original equation, after these substitutions, is
$$
c^n\left((a')^n+\sum_{i=0}^{n-1}r_i(a')^i(b')^{n-i}\right)=0.
$$
The nonzero factor $c^n$ in this domain makes the parenthesized
expression zero. Expanding the original identity
$1=(ua'+vb')^n$ and substituting that expression gives $1=b'w$,
where the entire inverse is
$$
w=-u^n\sum_{i=0}^{n-1}r_i(a')^i(b')^{n-i-1}
+\sum_{j=0}^{n-1}\binom nj(ua')^jv^{n-j}(b')^{n-j-1}\in R.
$$
Indeed the omitted top term of the latter binomial sum is
$u^n(a')^n$, exactly the term replaced using the monic equation.
Returning to the original numerator and denominator gives
$b(a'w)=cb'a'w=ca'=a$. Hence $g=a/b=a'w\in R$.
Every factor from the original fraction and every coefficient
of its monic equation is retained in this comparison.
\end{proof}

\begin{lemma}
\label{editorial-source-lemma-prepare-polynomial-ring-normal}
Let $R$ be a domain with fraction field $K$.
Suppose $f = \sum \alpha_i x^i$ is an
element of $K[x]$.
\begin{enumerate}
\item $f$ is integral over $R[x]$ if and only if
all $\alpha_i$ are integral over $R$, and
\item $f$ is almost integral over $R[x]$ if and only if
all $\alpha_i$ are almost integral over $R$.
\end{enumerate}
\end{lemma}

\begin{proof}
If $f=0$, all its coefficients are zero and the assertions hold
directly. Suppose $f\ne0$. We first prove the forward implication
in the second statement, keeping all coefficients, including zeros.
Write $f = \alpha_0 + \alpha_1 x + \ldots + \alpha_r x^r$
with $\alpha_r \not = 0$.
By assumption there exists $h = b_0 + b_1 x + \ldots + b_s x^s \in R[x]$,
$b_s \not = 0$ such that $f^n h \in R[x]$ for all
$n \geq 0$. This implies that $b_s \alpha_r^n \in R$
for all $n \geq 0$. Hence $\alpha_r$ is almost
integral over $R$. Since the set of almost integral
elements forms a subring (Lemma \ref{editorial-source-lemma-almost-integral}) we deduce that
$f - \alpha_r x^r = \alpha_0 + \alpha_1 x + \ldots + \alpha_{r - 1} x^{r - 1}$
is almost integral over $R[x]$. Here $\alpha_rx^r$ has the
actual witness $b_s$, since
$b_s(\alpha_rx^r)^n=(b_s\alpha_r^n)x^{rn}\in R[x]$.
The product $b_sh$ is consequently a witness for their difference
by the binomial calculation in Lemma \ref{editorial-source-lemma-almost-integral}.
The induction can be made fully explicit. For $0\leq t\leq r$
keep $f_t=\sum_{i=0}^t\alpha_ix^i$ and set
$$
h_t=b_s^{2^{r-t}-1}h,\qquad c_t=b_s^{2^{r-t}}.
$$
The original hypothesis says $h_rf_r^n\in R[x]$ for all $n\geq0$.
If $h_tf_t^n\in R[x]$, its coefficient of $x^{tn+s}$ is
$c_t\alpha_t^n$, also when $\alpha_t=0$. Hence this coefficient
belongs to $R$. The full binomial formula for
$f_{t-1}=f_t-\alpha_tx^t$ shows
$$
h_tc_t f_{t-1}^n
=\sum_{j=0}^n\binom nj(h_tf_t^j)
 \bigl(c_t(-\alpha_tx^t)^{n-j}\bigr)\in R[x].
$$
Since $h_tc_t=h_{t-1}$, this proves the descending induction.
In particular each original $\alpha_t$ has nonzero witness $c_t$,
and the single nonzero element $b_s^{2^r}$ is a witness for
all coefficients, by multiplying their equations by the remaining
power of $b_s$.

Conversely, if $\alpha_i$ has witness $0\ne w_i\in R$ for
$0\leq i\leq r$, put $w=\prod_iw_i\ne0$. For every $n\geq0$
the complete multinomial expansion gives
$$
w f^n=\sum_{k_0+\cdots+k_r=n}
\binom{n}{k_0,\ldots,k_r}
\left(\prod_{i=0}^r w_i\alpha_i^{k_i}\right)
x^{\sum_{i=0}^r i k_i}\in R[x].
$$
The indices $k_i$ are nonnegative integers and the multinomial
coefficients are integers, not inverses of factorials in $R$.
This proves the reverse almost-integral implication.

\medskip\noindent
In order to prove the first statement we will use absolute Noetherian
reduction. Namely, write $\alpha_i = a_i / b_i$ and
let $P(t) = t^d + \sum_{j < d} f_j t^j$ be a polynomial
with coefficients $f_j \in R[x]$ such that $P(f) = 0$.
Let $f_j = \sum f_{ji}x^i$. Consider the subring
$R_0 \subset R$ generated by the finite list of elements
$a_i, b_i, f_{ji}$ of $R$. It is a domain; let
$K_0$ be its field of fractions. Since $R_0$ is a finite type
$\mathbf{Z}$-algebra it is Noetherian, see
Lemma \ref{editorial-source-lemma-obvious-Noetherian}. It is still
the case that $f \in K_0[x]$ is integral over $R_0[x]$,
because the inclusion $R_0\subset R$ induces the injective maps
$K_0\to K$, $a/b\mapsto a/b$, and $K_0[x]\to K[x]$.
Each original denominator $b_i$ is nonzero in $R_0$, all the
original coefficients $f_{ji}$ belong to it, and the same monic
polynomial $P$ and same $f$ have $P(f)\in K_0[x]$.
Its image is zero in $K[x]$, so injectivity gives $P(f)=0$
already in $K_0[x]$.
By Lemma \ref{editorial-source-lemma-almost-integral} the element
$f$ is almost integral over $R_0[x]$. By the second statement of
the lemma, the elements $\alpha_i$ are almost integral
over $R_0$. And since $R_0$ is Noetherian, they are
integral over $R_0$, see Lemma \ref{editorial-source-lemma-almost-integral}.
Each resulting monic equation over $R_0$ maps coefficient by
coefficient to a monic equation for the same $\alpha_i$ over $R$.
For the reverse implication, if all the finitely many $\alpha_i$
are integral over $R$, the original algebra
$B=R[\alpha_0,\ldots,\alpha_r]\subset K$ is finite over $R$
by Lemma \ref{editorial-source-lemma-characterize-integral}. A finite module
generating family of $B$ also generates $B[x]$ over $R[x]$:
express each coefficient of a polynomial in that family.
Thus the original $f\in B[x]$ is integral over $R[x]$ by
Lemma \ref{algebra-lemma-finite-is-integral}, proving the other implication.
\end{proof}

\begin{lemma}
\label{editorial-source-lemma-polynomial-domain-normal}
Let $R$ be a normal domain.
Then $R[x]$ is a normal domain. For any domain $R$, the polynomial
ring $R[x]$ is completely normal if and only if $R$ is completely normal.
\end{lemma}

\begin{proof}
The result is true if $R$ is a field $K$ because
$K[x]$ is a euclidean domain and hence a principal ideal
domain and hence normal by Lemma \ref{editorial-source-lemma-PID-normal}.
The inclusions $R[x]\subset K[x]$ identify their fraction fields
with $K(x)$. More explicitly a fraction of polynomials over $R$
maps to the same fraction over $K$. Conversely, for $A/B$ with
$A,B\in K[x]$, choose a common nonzero denominator $c\in R$
of every coefficient of both polynomials. Its inverse image is
$(cA)/(cB)$, with both original scaled polynomials in $R[x]$.
The field equality criterion verifies independence of the chosen
common denominator and both composites.
Let $g$ be an element of the fraction field of
$R[x]$ which is integral over $R[x]$. Because $g$
is integral over $K[x]$ where $K$ is the fraction
field of $R$ we may write $g = \alpha_d x^d + \alpha_{d-1}x^{d-1} +
\ldots + \alpha_0$ with $\alpha_i \in K$.
By Lemma \ref{editorial-source-lemma-prepare-polynomial-ring-normal}
the elements $\alpha_i$ are integral over $R$ and
hence are in $R$. Thus the original $g$ belongs to $R[x]$.

For the second assertion, suppose $R$ is completely normal and
$g\in K(x)$ is almost integral over $R[x]$. Its same nonzero
polynomial witness also makes it almost integral over $K[x]$.
The ring $K[x]$ is Noetherian, since it is a principal ideal
domain, so Lemma \ref{editorial-source-lemma-almost-integral} makes $g$ integral
over $K[x]$. Normality of $K[x]$ then gives $g\in K[x]$.
The proved coefficientwise almost-integral equivalence shows
that every coefficient of this polynomial is almost integral
over $R$, hence lies in $R$. Therefore $g\in R[x]$.
Conversely, an almost integral element $g\in K$ and its witness
$0\ne r\in R$ remain the same constant fraction and witness
over $R[x]$. Complete normality of $R[x]$ places $g$ in
$R[x]\cap K=R$, where equality follows from uniqueness of
polynomial coefficients. This proves complete normality of $R$.
\end{proof}

\begin{lemma}
\label{editorial-source-lemma-power-series-over-Noetherian-normal-domain}
Let $R$ be a Noetherian normal domain. Then $R[[x]]$ is
a Noetherian normal domain. More generally, for any domain $R$,
$R[[x]]$ is completely normal if and only if $R$ is completely normal.
\end{lemma}

\begin{proof}
First prove the second assertion in the forward direction.
Assume that $R$ is completely normal, and retain nonzero
$f,g\in R[[x]]$ and the original $w=f/g$ in its fraction field,
now assumed almost integral. The power series ring is a domain:
the first nonzero coefficients of two nonzero series multiply
to the first nonzero coefficient of their product. Coefficient
inclusion gives an injection $R[[x]]\to K[[x]]$, where
$K=\operatorname{Frac}(R)$, hence an injection of its fraction
field into $K((x))$. A nonzero series over $K$ is $x^a$ times
a series with nonzero constant coefficient. For a series
$x^a\sum_{j\geq0}c_jx^j$ with $c_0\ne0$, its inverse is
$x^{-a}\sum_{j\geq0}d_jx^j$, where
$d_0=c_0^{-1}$ and
$d_j=-c_0^{-1}\sum_{i=1}^j c_id_{j-i}$ for $j\geq1$.
The coefficient of degree zero in the product is $1$ and every
higher coefficient is zero by this recurrence, proving the
inverse formula in $K((x))$. Thus the
original $w$ has the unique Laurent expansion
$$
w=a_nx^n+a_{n+1}x^{n+1}+\cdots,
\qquad n\in\mathbf Z,\quad a_n\ne0.
$$
Keep its original nonzero almost-integrality witness
$$
h=b_mx^m+b_{m+1}x^{m+1}+\cdots\in R[[x]],
\qquad b_m\ne0,
$$
so $hw^e\in R[[x]]$ for every $e\geq0$.
The first nonzero term of $hw^e$ has degree $m+en$ and
coefficient $b_ma_n^e$. If $n<0$, choose $e$ with $m+en<0$;
that nonzero negative-degree term contradicts membership in
$R[[x]]$. Hence $n\geq0$. For every $e\geq0$ the displayed
coefficient belongs to $R$, so complete normality gives $a_n\in R$.

Suppose all original coefficients $a_n,\ldots,a_{N-1}$ have
been proved to belong to $R$, retaining zeros in this list, and
put $P_N=\sum_{i=n}^{N-1}a_ix^i\in R[x]$.
The same original $h$ is a witness for $w-P_N$, since the full
identity
$$
h(w-P_N)^e=\sum_{j=0}^e\binom ej(hw^j)(-P_N)^{e-j}
$$
has every term in $R[[x]]$. The coefficient of $x^{m+eN}$
on the left is $b_ma_N^e$, including the case $a_N=0$:
$w-P_N$ has no terms of degree below $N$, so only these leading
terms can contribute at that degree. Thus $b_ma_N^e\in R$
for every $e\geq0$. The original nonzero $b_m$ witnesses
almost integrality of $a_N$ over $R$, giving $a_N\in R$.
Induction proves every coefficient belongs to $R$, hence the
original fraction $w$ belongs to $R[[x]]$. The zero fraction
belongs to the ring directly. This proves complete normality
of $R[[x]]$ without any Noetherian assumption.

Conversely, if $R[[x]]$ is completely normal and $z\in K$ is
almost integral over $R$, its original constant numerator,
denominator and witness define the same almost integral constant
fraction in the fraction field of $R[[x]]$. Thus it lies in
$R[[x]]\cap K=R$: uniqueness of the Laurent expansion forces
all coefficients but the constant one to be zero. Hence $R$
is completely normal, proving the equivalence.

Finally if the original $R$ is Noetherian and normal, Lemma
\ref{editorial-source-lemma-almost-integral} says that it is completely normal.
The argument just given makes $R[[x]]$ completely normal, hence
normal since every integral element is almost integral.
Its Noetherian property follows from
Lemma \ref{editorial-source-lemma-Noetherian-power-series}, proving the first assertion.
\end{proof}

\begin{lemma}
\label{editorial-source-lemma-normality-is-local}
Let $R$ be a domain. The following are equivalent:
\begin{enumerate}
\item The domain $R$ is a normal domain,
\item for every prime $\mathfrak p \subset R$ the local ring
$R_{\mathfrak p}$ is a normal domain, and
\item for every maximal ideal $\mathfrak m$ the ring $R_{\mathfrak m}$
is a normal domain.
\end{enumerate}
\end{lemma}

\begin{proof}
We deduce (1) $\Rightarrow$ (2) from Lemma \ref{editorial-source-lemma-localize-normal-domain}.
The implication (2) $\Rightarrow$ (3) is immediate. The implication
(3) $\Rightarrow$ (1) follows from the fact that for any domain $R$ we have
$$
R = \bigcap\nolimits_{\mathfrak m} R_{\mathfrak m}
$$
inside the original fraction field of $R$. To prove this equality,
let $g$ belong to every $R_{\mathfrak m}$ and retain
$I=\{x\in R:xg\in R\}$. It is an ideal: multiplication by
$g$ respects sums and multiplication by elements of $R$.
For each maximal $\mathfrak m$, the expression $g=a/s$ in
$R_{\mathfrak m}\subset\operatorname{Frac}(R)$ has
$s\notin\mathfrak m$ and gives $sg=a\in R$. Thus $s\in I$
and $I\not\subseteq\mathfrak m$. If $I$ were proper it would
lie in a maximal ideal, a contradiction. Therefore $1\in I$
and $g\in R$. Finally if $g$ is integral over $R$, its original
monic equation remains monic over each $R_{\mathfrak m}$.
Under (3), each of those rings contains $g$, and the proved
intersection identity gives $g\in R$, establishing (1).
\end{proof}

\noindent
Lemma \ref{editorial-source-lemma-normality-is-local} shows that the following definition
is compatible with Definition \ref{editorial-source-definition-domain-normal}. (It is the
definition from EGA -- see \cite[IV, 5.13.5 and 0, 4.1.4]{EGA}.)

\begin{definition}
\label{editorial-source-definition-ring-normal}
A ring $R$ is called {\it normal} if for every prime
$\mathfrak p \subset R$ the localization $R_{\mathfrak p}$ is
a normal domain (see Definition \ref{editorial-source-definition-domain-normal}).
\end{definition}

\noindent
It suffices in this definition to test maximal ideals. Indeed,
if each $R_{\mathfrak m}$ is a normal domain and
$\mathfrak p\subseteq\mathfrak m$, there is the exact isomorphism
$$
(R_{\mathfrak m})_{\mathfrak pR_{\mathfrak m}}
\longrightarrow R_{\mathfrak p},\qquad
\frac{a/s}{b/t}\longmapsto\frac{at}{sb},
$$
with inverse $a/u\mapsto(a/1)/(u/1)$.
Here $s,t\notin\mathfrak m$ and $b\notin\mathfrak p$, so
$sb\notin\mathfrak p$. The prime correspondence for localization
identifies $\mathfrak pR_{\mathfrak m}$ as a proper prime with
contraction $\mathfrak p$. Fraction equality witnesses map to
fraction equality witnesses in these localizations; the displayed
formulas fix every original fraction in both composites.
Lemma \ref{editorial-source-lemma-localize-normal-domain} then makes
$R_{\mathfrak p}$ a normal domain. The converse is part of the
definition, and every prime lies in a maximal ideal.

A normal ring is reduced. More explicitly the natural map
$R\to\prod_{\mathfrak m}R_{\mathfrak m}$ is injective: if
$a\ne0$, its annihilator is a proper ideal and is contained in
a maximal ideal $\mathfrak m$; then $a/1$ cannot be zero there,
since a denominator outside $\mathfrak m$ killing $a$ would
belong to that annihilator. A nilpotent element maps to zero
in every domain $R_{\mathfrak m}$ and hence is zero in $R$.
This also proves the stated subring assertion using all primes
(see Lemma \ref{algebra-lemma-characterize-zero-local}).
For the zero ring the products are empty and all assertions
hold directly.

\begin{lemma}
\label{editorial-source-lemma-normal-ring-integrally-closed}
A normal ring is integrally closed in its total ring of fractions.
\end{lemma}

\begin{proof}
Let $R$ be normal, let $\Sigma\subset R$ be its original
multiplicative set of nonzerodivisors, and keep
$Q(R)=\Sigma^{-1}R$. The map $R\to Q(R)$ is injective,
since any element in its kernel is killed by a member of $\Sigma$.
Let $x\in Q(R)$ be integral over $R$, and retain the ideal
$I=\{f\in R\mid fx\in R\}$.
Fix a prime $\mathfrak p$ and put $A=R_{\mathfrak p}$, a
normal domain. Every $s\in\Sigma$ has nonzero image in $A$:
an equality $s/1=0$ would give $u\notin\mathfrak p$ with
$us=0$, forcing $u=0$, a contradiction.
There is the exact comparison
$$
Q(R)\otimes_RA\cong\Sigma^{-1}A
\lhook\joinrel\longrightarrow\operatorname{Frac}(A).
$$
The first map sends $(r/s)\otimes c$ to $(rc)/(s/1)$;
its inverse sends $c/(s/1)$ to $(1/s)\otimes c$.
The original coefficient maps agree on $R$, so tensor balancing
gives the first map. The second is defined because each $s/1$
maps to the unit $(s/1)\otimes1$ with inverse $(1/s)\otimes1$.
The formulas verify both composites. The last map is injective
since $A$ is a domain and every image of $\Sigma$ is nonzero.
For the original fractions its full formula is
$$
(r/s)\otimes(a/u)\longmapsto
\frac{(ra)/1}{(su)/1},\qquad
s\in\Sigma,\quad u\notin\mathfrak p.
$$
This embeds the indicated localization in the fraction field;
it does not require that every nonzero element of $A$ come from
an original nonzerodivisor of $R$.

The image of $x\otimes1$ satisfies the same monic equation
with coefficients mapped into $A$. Normality of $A$ places
it in $A$. The injective comparison just proved therefore gives
$x\otimes1=a\otimes(1/f)$ in
$Q(R)\otimes_RR_{\mathfrak p}\cong Q(R)_{\mathfrak p}$ for
some original $a,f\in R$, $f\notin\mathfrak p$.
This last isomorphism sends $z\otimes(a/u)$ to $az/u$
and has inverse $z/u\mapsto z\otimes(1/u)$; tensor balancing
and the fraction equality relation verify both maps and composites.
Hence $fx-a$ becomes zero in that localization, so there is
$f'\notin\mathfrak p$ with $f'(fx-a)=0$ in $Q(R)$.
The equality $f'fx=f'a$ has its right side in the embedded $R$,
so $ff'\in I$ and $ff'\notin\mathfrak p$.
As this holds for every prime, $I$ cannot be proper; otherwise
a maximal ideal containing it would contradict this property.
Thus $1\in I$ and the original $x$ lies in $R$.
If $R=0$, its total quotient ring is zero and the assertion
holds without any choice of a prime.
\end{proof}

\begin{lemma}
\label{editorial-source-lemma-localization-normal-ring}
A localization of a normal ring is a normal ring.
\end{lemma}

\begin{proof}
Let $S\subset R$ be the original multiplicative subset and
let $\mathfrak q$ be a prime of $S^{-1}R$. Its contraction
$\mathfrak p$ to $R$ is prime and disjoint from $S$, and
$\mathfrak q=S^{-1}\mathfrak p$. The maps
$$
(S^{-1}R)_{\mathfrak q}\longrightarrow R_{\mathfrak p},
\quad\frac{a/s}{b/t}\longmapsto\frac{at}{sb},
\qquad
R_{\mathfrak p}\longrightarrow(S^{-1}R)_{\mathfrak q},
\quad a/u\longmapsto\frac{a/1}{u/1}
$$
are defined because $s,t\in S$ and $b\notin\mathfrak p$,
while $u\notin\mathfrak p$ in the second formula. The original
maps from $R$ send each required denominator to a unit, hence
extend to these localizations; on fractions the displayed formulas
prove that the two extensions are inverse. Normality of $R$
makes $R_{\mathfrak p}$ a normal domain, proving the assertion
at every prime of $S^{-1}R$. If this localization is zero,
there are no primes to test and it is normal by definition.
\end{proof}

\begin{lemma}
\label{editorial-source-lemma-polynomial-ring-normal}
Let $R$ be a normal ring. Then $R[x]$ is a normal ring.
More generally, the polynomial ring over $R$ in any set of
indeterminates is normal.
\end{lemma}

\begin{proof}
Let $\mathfrak q\subset R[x]$ be prime and let
$\mathfrak p$ be its contraction to the original $R$.
The ring $R_{\mathfrak p}[x]$ is a normal domain by
Lemma \ref{editorial-source-lemma-polynomial-domain-normal}. It is canonically
$(R\setminus\mathfrak p)^{-1}R[x]$: a finite polynomial with
fraction coefficients is represented by multiplying all their
original denominators to obtain one common denominator, and
the inverse sends $F/s$ to the polynomial whose coefficients
are the corresponding fractions. These maps preserve coefficient
equalities and are inverse. The extended ideal
$\widetilde{\mathfrak q}=(R\setminus\mathfrak p)^{-1}\mathfrak q$
is prime, with contraction $\mathfrak q$. The exact local maps are
$$
(R[x])_{\mathfrak q}\longrightarrow
(R_{\mathfrak p}[x])_{\widetilde{\mathfrak q}},
\quad F/G\longmapsto(F/1)/(G/1),
$$
with inverse $(F/s)/(G/t)\mapsto tF/(sG)$.
Here $s,t\in R\setminus\mathfrak p$ and $G\notin\mathfrak q$,
so every displayed denominator is valid. The original coefficient
maps invert the indicated sets, proving well-definedness by the
localization relations, and the fraction formulas verify both
composites. Lemma \ref{editorial-source-lemma-localize-normal-domain} makes this
local ring a normal domain, proving the one-variable assertion.

Induction proves the assertion for any finite set of variables.
For an arbitrary set $J$, retain all its original labeled variables.
The polynomial ring $R[x_j:j\in J]$ is the directed colimit of
$R[x_j:j\in F]$ over the finite subsets $F\subset J$:
every polynomial has finite support, and coefficient equality
already holds in a finite set containing the variables of both
polynomials. The transition maps keep each original variable and
coefficient. Lemma \ref{editorial-source-lemma-colimit-normal-ring} applies to
these normal rings and gives the full assertion. Its proof below
does not use this polynomial-ring assertion.
\end{proof}

\begin{lemma}
\label{editorial-source-lemma-finite-product-normal}
A finite product of normal rings is normal. More generally,
for any set $I$, the product $\prod_{i\in I}R_i$ is normal
if and only if every original factor $R_i$ is normal.
Every ultraproduct of normal rings is normal as well.
\end{lemma}

\begin{proof}
It suffices to show that the product of two normal rings, say $R$ and $S$, is
normal. By Lemma \ref{algebra-lemma-disjoint-decomposition} the prime ideals of
$R\times S$ are of the form $\mathfrak{p}\times S$ and $R\times
\mathfrak{q}$, where $\mathfrak{p}$ and $\mathfrak{q}$ are primes of $R$
and $S$ respectively. Localization yields 
$(R\times S)_{\mathfrak{p}\times S}=R_{\mathfrak{p}}$ which is a normal domain
by assumption. Explicitly the map sends
$(a,b)/(s,t)$ to $a/s$, and its inverse sends $a/s$ to
$(a,0)/(s,1)$. Both are defined because $s\notin\mathfrak p$.
The idempotent $(1,0)$ is a unit in this localization, so $(0,1)$
is zero there, verifying the inverse identities. The corresponding
maps at $R\times\mathfrak q$ use the other coordinate, proving
the finite assertion without any assumption on the factor maps.
\medskip\noindent
For the arbitrary-product assertion we first prove an algebraic
characterization. A ring $B$ is normal if and only if, for every
$n\geq1$ and all $a,b,c_0,\ldots,c_{n-1}\in B$, it satisfies
$$
a^n+\sum_{j=0}^{n-1}c_ja^jb^{n-j}=0
\quad\Longrightarrow\quad a=bd\text{ for some }d\in B.
\qquad(\mathcal D_n)
$$
Here $b$ may be zero or a zero divisor.
Suppose first that $B$ is normal and the displayed equation holds.
At a prime $\mathfrak p$, write $a_{\mathfrak p},b_{\mathfrak p}$
for the original images in $B_{\mathfrak p}$. If
$b_{\mathfrak p}=0$, the equation gives $a_{\mathfrak p}^n=0$,
hence $a_{\mathfrak p}=0$ since this localization is a domain.
If $b_{\mathfrak p}\ne0$, division by $b_{\mathfrak p}^n$
gives the monic equation
$$
(a_{\mathfrak p}/b_{\mathfrak p})^n+
\sum_{j=0}^{n-1}c_{j,\mathfrak p}
 (a_{\mathfrak p}/b_{\mathfrak p})^j=0.
$$
Integral closedness then gives
$a_{\mathfrak p}\in b_{\mathfrak p}B_{\mathfrak p}$ in both cases.
This implies $a\in bB$: otherwise the ideal
$J=\{r\in B:ra\in bB\}$ is proper and lies in a maximal ideal
$\mathfrak p$. Membership in the localized principal ideal gives
$s\notin\mathfrak p$ with $sa\in bB$ (clear the denominator and
retain the further multiplier witnessing equality in the localization).
This contradicts $s\in J\subset\mathfrak p$ and proves
$(\mathcal D_n)$.

Conversely suppose all $(\mathcal D_n)$ hold. Condition
$(\mathcal D_2)$ with $b=0$ and both coefficients zero gives
$z^2=0\Rightarrow z=0$. If $ab=0$, apply $(\mathcal D_2)$
to the exact equation $a^2-a(a+b)=0$ to get
$a=(a+b)d$ for some $d\in B$. Put
$z=(1-d)a=db$. Then
$z^2=((1-d)a)(db)=d(1-d)ab=0$, so $z=0$ and
$$
(1-d)a=0,\qquad db=0,\qquad (1-d)+d=1.
$$
These identities show that every $B_{\mathfrak p}$ is a domain.
Indeed a zero product $(a/u)(b/v)=0$ has a witness
$w\notin\mathfrak p$ with $wab=0$ in $B$. Apply the identities
to $wa$ and $b$. At least one of $d,1-d$ is outside
$\mathfrak p$. If $d$ is outside, $b/v=0$ in the localization;
if $1-d$ is outside, $a/u=0$ since $w$ is also a unit there.
The localization is nonzero because $\mathfrak p$ is proper.

To prove this domain integrally closed, retain a fraction
$y=(a_0/s)/(b_0/t)$ with $s,t\notin\mathfrak p$ and
$b_0/t\ne0$. Set $a=a_0t$ and $b=b_0s$ in the original $B$,
so $b/1\ne0$ and $y=(a/1)/(b/1)$ with these factors recorded.
Suppose its original monic equation is
$$
y^n+\sum_{j=0}^{n-1}(r_j/t_j)y^j=0,
\qquad r_j\in B,\quad t_j\notin\mathfrak p.
$$
Put $u=\prod_{j=0}^{n-1}t_j$ and
$c_j=r_j\prod_{k\ne j}t_k$, retaining every denominator factor.
Then $r_j/t_j=c_j/u$ in the localization. Multiplication by
$(b/1)^n(u/1)^n$ says that the element
$$
F=(ua)^n+\sum_{j=0}^{n-1}c_j u^{n-1-j}(ua)^jb^{n-j}
$$
has zero image in $B_{\mathfrak p}$. Therefore some
$v\notin\mathfrak p$ satisfies $vF=0$ in the original ring.
The full equation
$$
(vua)^n+\sum_{j=0}^{n-1}c_j u^{n-1-j}(vua)^j(vb)^{n-j}
=v^nF=0
$$
holds there, because $n\geq1$. Condition $(\mathcal D_n)$ gives
$vua=vbd$ for some $d\in B$. Since $v/1$ and $u/1$ are units
and $b/1\ne0$ in the localized domain, the original fraction
satisfies $y=d/u\in B_{\mathfrak p}$. This proves the characterization,
including the zero ring, whose prime-local condition is vacuous.

Now let $R=\prod_{i\in I}R_i$ with every $R_i$ normal.
An equation as in $(\mathcal D_n)$ in $R$ means precisely
$$
a_i^n+\sum_{j=0}^{n-1}c_{j,i}a_i^jb_i^{n-j}=0
\quad\text{in every original }R_i.
$$
Choose $d_i\in R_i$ with $a_i=b_id_i$ using the proved
characterization, and put $d=(d_i)_{i\in I}$. It gives $a=bd$
in the original product. Thus $R$ satisfies every $(\mathcal D_n)$
and is normal. Only the equation degree and its finitely many
coefficients are common; no characteristic or finiteness condition
on $I$ or on the factors is used.

We describe the exact prime and localization maps as well,
including the primes of an infinite product that are not obtained
from a single coordinate. For $E\subseteq I$, let $e_E\in R$
have coordinate $1$ on $E$ and $0$ on its complement. For a
prime $\mathfrak p\subset R$, the family
$$
\mathcal U=\{E\subseteq I:e_E\notin\mathfrak p\}
$$
is an ultrafilter, that is, a family closed under finite intersections
and enlargement, containing $I$ but not the empty set, and containing
exactly one of each subset and its complement. These facts follow
respectively from
$$
e_Ee_F=e_{E\cap F},\quad
e_E=e_Ee_F\ (E\subseteq F),\quad
e_I=1,\quad e_\varnothing=0,
$$
and from $e_Ee_{I\setminus E}=0$,
$e_E+e_{I\setminus E}=1$, using primality.
Set
$$
K_{\mathcal U}=\{a=(a_i):\{i:a_i=0\}\in\mathcal U\},
\qquad A_{\mathcal U}=R/K_{\mathcal U},
\qquad\pi:R\to A_{\mathcal U}.
$$
The set $K_{\mathcal U}$ is an ideal: intersections of zero sets
control sums, and multiplication by any tuple enlarges a zero set.
The quotient $A_{\mathcal U}$ is the ultraproduct, here defined by
this exact ideal. Also $K_{\mathcal U}\subseteq\mathfrak p$:
if $E=\{i:a_i=0\}\in\mathcal U$, then $e_Ea=0$ and
$e_E\notin\mathfrak p$, so primality gives $a\in\mathfrak p$.
Therefore $\mathfrak q=\mathfrak p/K_{\mathcal U}$ is a prime
of $A_{\mathcal U}$, and the map
$$
R_{\mathfrak p}\longrightarrow (A_{\mathcal U})_{\mathfrak q},
\qquad a/s\longmapsto\pi(a)/\pi(s)
$$
is defined, because $s\notin\mathfrak p$ if and only if
$\pi(s)\notin\mathfrak q$. It is surjective by lifting each
numerator and denominator through the quotient. It is injective:
if $\pi(a)/\pi(s)=0$, choose $t\notin\mathfrak p$ with
$ta\in K_{\mathcal U}$. For
$E=\{i:t_ia_i=0\}\in\mathcal U$, one has
$e_Eta=0$ and $e_Et\notin\mathfrak p$. This is the original
localization witness that $a/s=0$. The inverse map is consequently
$\pi(a)/\pi(s)\mapsto a/s$, independent of the chosen lifts by
the same argument applied to their difference. Both composites fix
the original fractions.

Conversely a prime of an ultraproduct contracts to a prime of
the product and recovers its ultrafilter: the class of $e_E$ is
$1$ when $E\in\mathcal U$ and $0$ when its complement is in
$\mathcal U$. Thus, as sets, all product primes are uniquely the
pairs of an ultrafilter and a prime of its ultraproduct, with the
displayed localization isomorphism for each pair.

The ultraproduct of normal rings is itself normal. If a displayed
$(\mathcal D_n)$ equation holds in $A_{\mathcal U}$, it holds
coordinatewise on a set $E\in\mathcal U$ by the definition of
the quotient ideal. For $i\in E$ choose $d_i$ with $a_i=b_id_i$
in $R_i$; set $d_i=0$ elsewhere. Their classes satisfy
$\pi(a)=\pi(b)\pi(d)$. The characterization proves normality
of $A_{\mathcal U}$, so the localization isomorphism also gives
the product result through its actual primes.

When every factor is a nonzero normal domain, these maps have
a stronger description. The ultraproduct is then a domain:
the zero set of a product is the union of the zero sets of its
two factors, so membership in an ultrafilter forces one of those
zero sets to belong to it. Its identity is nonzero because the
coordinate identities are nonzero. If $\pi(a)/\pi(b)$ is integral
in its fraction field with $\pi(b)\ne0$, then the sets
$$
T=\{i:b_i\ne0\},\qquad
E=\{i:a_i^n+\sum_{j=0}^{n-1}c_{j,i}a_i^jb_i^{n-j}=0\}
$$
belong to $\mathcal U$. On $E\cap T$, the original fraction
$a_i/b_i$ satisfies its monic equation over $R_i$ and belongs
to $R_i$. Choose that value for $d_i$ there and zero elsewhere.
Then $\pi(a)=\pi(b)\pi(d)$, proving integral closedness directly.
Moreover in this domain-factor case
$\ker(R\to R_{\mathfrak p})=K_{\mathcal U}$.
One inclusion was proved by the idempotent witnesses above.
For the other, if $ta=0$ with $t\notin\mathfrak p$, then
$t\notin K_{\mathcal U}$, so $\{i:t_i\ne0\}\in\mathcal U$;
on that set the domains force $a_i=0$, giving $a\in K_{\mathcal U}$.

Finally normality of the product implies normality of each factor.
For $\mathfrak r\in\operatorname{Spec}(R_i)$, take
$\mathfrak p=\operatorname{pr}_i^{-1}(\mathfrak r)$.
The localization map
$$
R_{\mathfrak p}\longrightarrow(R_i)_{\mathfrak r},
\qquad a/s\longmapsto a_i/s_i
$$
has inverse obtained by using the numerator tuple supported at $i$
and the denominator tuple equal to $s_i$ at $i$ and to $1$ elsewhere.
Multiplication by $e_{\{i\}}\notin\mathfrak p$ kills any
differences outside that coordinate, which verifies both composites
and independence of fraction representatives. Hence these local
rings are isomorphic, proving normality of the factor.
Zero-ring factors cause no exception: if $J=\{i:R_i\ne0\}$,
then $e_J=1_R$, so any ultrafilter obtained from a product prime
contains $J$. An ultrafilter concentrated on zero factors gives
the zero ultraproduct and contributes no primes. The empty product
is the zero ring and satisfies the same conclusion.
\end{proof}

\begin{lemma}
\label{editorial-source-lemma-characterize-reduced-ring-normal}
Let $R$ be a ring. Assume $R$ is reduced and has finitely many
minimal primes. Then the following are equivalent:
\begin{enumerate}
\item $R$ is a normal ring,
\item $R$ is integrally closed in its total ring of fractions, and
\item $R$ is a finite product of normal domains.
\end{enumerate}
\end{lemma}

\begin{proof}
The implications (1) $\Rightarrow$ (2) and
(3) $\Rightarrow$ (1) hold in general,
see Lemmas \ref{editorial-source-lemma-normal-ring-integrally-closed} and
\ref{editorial-source-lemma-finite-product-normal}.

\medskip\noindent
Let $\mathfrak p_1, \ldots, \mathfrak p_n$ be the minimal primes of $R$.
By Lemmas \ref{algebra-lemma-reduced-ring-sub-product-fields} and
\ref{algebra-lemma-total-ring-fractions-no-embedded-points} we have
$Q(R) = R_{\mathfrak p_1} \times \ldots \times R_{\mathfrak p_n}$, and
by Lemma \ref{algebra-lemma-minimal-prime-reduced-ring} each factor is a field.
Let $e_i = (0, \ldots, 0, 1, 0, \ldots, 0)$ be the $i$th idempotent
of $Q(R)$.

\medskip\noindent
Suppose $R$ is integrally closed in its original $Q(R)$.
Each $e_i$ satisfies $e_i^2-e_i=0$, hence lies in $R$.
The orthogonality identities and $\sum_i e_i=1$ give inverse maps
$$
R\longrightarrow\prod_{i=1}^n e_iR,\quad
r\longmapsto(e_ir)_i,
\qquad (r_i)_i\longmapsto\sum_{i=1}^n r_i.
$$
The factor $e_iR$ has identity $e_i$ and embeds in the original
field $R_{\mathfrak p_i}$ by the $i$th coordinate of $Q(R)$.
The kernel of $R\to e_iR$ is exactly $\mathfrak p_i$.
Indeed the kernel of $R\to R_{\mathfrak p_i}$ is
$\mathfrak p_i$: the localization is a field, whose maximal ideal
$\mathfrak p_iR_{\mathfrak p_i}$ is zero; conversely an equation
$ur=0$ with $u\notin\mathfrak p_i$ forces $r\in\mathfrak p_i$
by primality. Multiplication by $e_i$ keeps precisely this
coordinate, and injectivity of $R\to Q(R)$ identifies its
vanishing with $e_ir=0$ in $R$. Therefore
$R/\mathfrak p_i\cong e_iR$ by $r+\mathfrak p_i\mapsto e_ir$,
with inverse $e_ir\mapsto r+\mathfrak p_i$.
The maps
$$
\operatorname{Frac}(R/\mathfrak p_i)\longrightarrow R_{\mathfrak p_i},
\quad\frac{a+\mathfrak p_i}{b+\mathfrak p_i}
\longmapsto\frac{a/1}{b/1},
$$
and $a/s\mapsto(a+\mathfrak p_i)/(s+\mathfrak p_i)$ are inverse:
the denominators avoid $\mathfrak p_i$, and the original fraction
equality witnesses give equality in the domain $R/\mathfrak p_i$.
This proves the asserted field-of-fractions comparison.

Each subring $R/\mathfrak p_i\subset R_{\mathfrak p_i}$ is
integrally closed, as also follows from
Lemma \ref{editorial-source-lemma-finite-product-integral-closure}.
Directly, if $z_i$ in that field satisfies a monic equation of
degree $d\geq1$ over $e_iR$, form $z\in Q(R)$ with $i$th
coordinate $z_i$ and all others zero. Regard each coefficient
as the element of $R$ supported by $e_i$. The same monic
polynomial vanishes on $z$: at the other coordinates it is
evaluated at zero with all coefficients zero, including its
constant coefficient. Thus $z$ is integral over $R$ and belongs
to $R$ by hypothesis. Its $i$th coordinate lies in $e_iR$,
so $z_i$ belongs to the stated subring. This proves that each
$R/\mathfrak p_i$ is a normal domain and gives (3).
If $R=0$, there are no minimal primes and (3) is the empty
product, so the same conclusion holds directly.
\end{proof}

\begin{lemma}
\label{editorial-source-lemma-colimit-normal-ring}
Let $(R_i,\varphi_{ii'})$ be a system of rings indexed by a nonempty
directed preordered set $I$
(Categories, Definition \ref{categories-definition-directed-system}).
If each $R_i$ is normal, so is
$R = \colim_i R_i$.
\end{lemma}

\begin{proof}
Write $\lambda_i:R_i\to R$ for the original colimit maps.
If $R=0$, the conclusion is vacuous. Otherwise fix a prime
$\mathfrak p\subset R$ and set
$\mathfrak p_i=\lambda_i^{-1}(\mathfrak p)$.
These are proper primes, and
$\varphi_{ij}^{-1}(\mathfrak p_j)=\mathfrak p_i$ because
$\lambda_j\varphi_{ij}=\lambda_i$.
Thus the localized transition maps are defined by
$$
\psi_{ij}:(R_i)_{\mathfrak p_i}\longrightarrow(R_j)_{\mathfrak p_j},
\qquad a_i/s_i\longmapsto
\varphi_{ij}(a_i)/\varphi_{ij}(s_i).
$$
Indeed denominators outside $\mathfrak p_i$ stay outside
$\mathfrak p_j$, and an equality witness
$u_i(a_it_i-b_is_i)=0$ with $u_i\notin\mathfrak p_i$
maps to the corresponding equality witness in $R_j$.
Put $A_i=(R_i)_{\mathfrak p_i}$ and $L=\colim_i A_i$,
with maps $\rho_i:A_i\to L$.

The element rule for a directed colimit says that every element
has a representative at one stage, and equality of two representatives
holds at a common later stage. To recall why, construct the underlying
set from pairs $(i,a_i)$, identifying two pairs when their images
agree at some common stage. Directedness makes this relation
transitive; addition and multiplication are computed at a common
stage and are independent of that stage by the same rule.
This construction satisfies the universal property of the ring
colimit, since every compatible family of maps sends equivalent
pairs to the same value. In particular a zero equality may need
a later stage; no transition map is assumed injective.

We prove the original localization comparison with both maps:
$$
\alpha:L\longrightarrow R_{\mathfrak p},\qquad
\rho_i(a_i/s_i)\longmapsto
\lambda_i(a_i)/\lambda_i(s_i),
$$
and
$$
\beta:R_{\mathfrak p}\longrightarrow L,\qquad
a/s\longmapsto\rho_i(a_i/s_i),
$$
where numerator and denominator are represented at one common
stage $i$. Since $s\notin\mathfrak p$, every such representative
$s_i$ is outside $\mathfrak p_i$. For well-definedness of the
second map, an equality $a/s=b/t$ has a witness
$u\notin\mathfrak p$ with $u(at-bs)=0$ in $R$.
Represent all five elements at one common stage. The equality
rule gives a later stage $j$ with
$u_j(a_jt_j-b_js_j)=0$ in $R_j$.
All three denominators $u_j,s_j,t_j$ are outside
$\mathfrak p_j$, so $a_j/s_j=b_j/t_j$ in $A_j$.
The same argument handles any change of representative.
Addition and multiplication use the ordinary fraction formulas
at a common stage. The first map is defined by the compatible
localized maps to $R_{\mathfrak p}$, whose fraction relations
are respected in exactly the same way.
The composites fix every $a/s$ and every generator
$\rho_i(a_i/s_i)$, respectively. Hence
$$
R_{\mathfrak p}\cong\colim_i(R_i)_{\mathfrak p_i}.
$$

Each $A_i$ is a nonzero normal domain by the hypothesis.
The colimit $L$ is nonzero: otherwise $1=0$ would hold at
some later stage, contradicting this property of $A_i$.
It is a domain even though the transition maps may have kernels.
For $xy=0$ in $L$, represent $x,y$ at one stage and then
pass to a later stage where their product is zero.
The domain at that stage forces one representative to be zero,
so one of the original $x,y$ is zero in $L$.
Thus $A=R_{\mathfrak p}$ is a nonzero domain.

Retain $a,b\in A$, $b\ne0$, and an original monic equation
$$
z^n+\sum_{k=0}^{n-1}c_kz^k=0,
\qquad z=a/b,\quad c_k\in A,\quad n\geq1.
$$
Multiplication by the full denominator power $b^n$ gives
$a^n+\sum_{k=0}^{n-1}c_ka^kb^{n-k}=0$ in $A$.
Choose a common stage $i$ representing each of the finitely
many localized elements as
$$
a=\frac{\lambda_i(a_i)}{\lambda_i(u_i)},\qquad
b=\frac{\lambda_i(b_i)}{\lambda_i(v_i)},\qquad
c_k=\frac{\lambda_i(c_{k,i})}{\lambda_i(w_{k,i})},
$$
with $u_i,v_i,w_{k,i}\notin\mathfrak p_i$.
Keep the products
$$
W_i=\prod_{k=0}^{n-1}w_{k,i},\qquad
W_{k,i}=\prod_{\ell\ne k}w_{\ell,i},\qquad
D_i=u_i^nv_i^nW_i
$$
and the full numerator
$$
H_i=a_i^nv_i^nW_i+
\sum_{k=0}^{n-1}c_{k,i}a_i^kb_i^{n-k}
 u_i^{n-k}v_i^kW_{k,i}.
$$
In $A_i$, without omitting any denominator factor, one has
$$
\frac{H_i}{D_i}
=\left(\frac{a_i}{u_i}\right)^n+
\sum_{k=0}^{n-1}\frac{c_{k,i}}{w_{k,i}}
 \left(\frac{a_i}{u_i}\right)^k
 \left(\frac{b_i}{v_i}\right)^{n-k}.
$$
The image of this fraction in $A$ is zero. Therefore some
$t\in R\setminus\mathfrak p$ satisfies
$t\lambda_i(H_i)=0$ in the original $R$.
Represent $t$ at a stage above $i$, then apply the equality rule
again to obtain a stage $q$ where $t_qH_q=0$ in $R_q$.
Here every entry of $H_q$ is the image of the corresponding
original entry of $H_i$ under the indicated transition map:
$$
H_q=a_q^nv_q^n\prod_{\ell=0}^{n-1}w_{\ell,q}
+\sum_{k=0}^{n-1}c_{k,q}a_q^kb_q^{n-k}u_q^{n-k}v_q^k
 \prod_{\ell\ne k}w_{\ell,q}.
$$
The witness $t_q$ and all the denominator factors are outside
$\mathfrak p_q$. Thus $t_q$ and $D_q$ are units in $A_q$,
and the complete displayed equation for $H_q/D_q$ is zero there.
The element $b_q/v_q$ is nonzero in $A_q$, since its image is
the original nonzero $b$ in $A$. Consequently
$$
z_q=\frac{a_q/u_q}{b_q/v_q}\in\operatorname{Frac}(A_q)
$$
is defined, and division by $(b_q/v_q)^n$ gives its monic equation
$$
z_q^n+\sum_{k=0}^{n-1}(c_{k,q}/w_{k,q})z_q^k=0.
$$
Normality gives $z_q=h_q\in A_q$. Transport the equality
$a_q/u_q=(b_q/v_q)h_q$ in $A_q$ through its actual map to $A$.
It becomes $a=bh$ for the image $h\in A$ of $h_q$.
Since $A$ is a domain and $b\ne0$, the original $z=a/b$ equals
$h$ and belongs to $A$. This proves integral closedness of $A$
and hence normality of $R$.
No map from $\operatorname{Frac}(A_q)$ to
$\operatorname{Frac}(A)$ has been assumed: it need not exist
when the stage map has a kernel. The equality transported is
the displayed equality inside $A_q$ itself.

The divisibility characterization in the arbitrary-product proof
gives a second direct verification. Represent the finitely many
elements of a $(\mathcal D_n)$ equation in $R$ at one stage,
then pass to a later stage where that whole equation holds.
The normal ring at that stage supplies $a_i=b_id_i$;
its image gives the required $a=bd$ in $R$. Thus all
$(\mathcal D_n)$ hold, in agreement with the localization proof.
Neither proof requires antisymmetry of the directed preorder
or injectivity of the transition maps.
\end{proof}







